\begin{proof}[正则场有限系数的证明]
取平直时空的局部解析点电荷模型，使用正文的推迟、超前场公式。目标是在半差中实际算出有限系数，而不是由最终 ALD 方程反推。\footnote{Dirac 的局部场展开及点粒子守恒方法见原始论文 \href{https://doi.org/10.1098/rspa.1938.0124}{Proc. R. Soc. A 167, 148--169 (1938)}；下列计算沿用正文的 SI 制与 $(+---)$ 度规。}
先确定局域有限项允许的协变结构。平直时空点 Maxwell 模型没有源的内禀长度或时间尺度。

对世界线作局域解析展开时，\(O(\ell^0)\) 的正则场必须是 Lorentz 协变、反对称且宇称偶的二阶张量。由 \(u^\mu\)、\(a^\mu\)、\(\dot a^\mu\) 构造并具有正确导数阶数的最低阶候选只有 \(u^{[\mu}\dot a^{\nu]}\)：\(a^{[\mu}a^{\nu]}=0\)，\(u^{[\mu}a^{\nu]}\) 少一个固有时导数而又没有可补偿的内禀时间尺度，含 Levi--Civita 张量的候选则是宇称奇的伪张量。这里要求局域表达在世界线 Taylor 数据处解析，排除 $\sqrt{-a^2}$ 等非解析结构。更高导数阶的解析项缺少补足维数的内禀长度尺度。

因此世界线上的有限正则场必可写成
\[
 F_{\rm R}^{\mu\nu}=C\,u^{[\mu}\dot a^{\nu]},
\]
其中唯一待定的是常数 \(C\)；而 \(u\cdot\dot a=-a^2\) 已自动包含 \(\dot a^\mu\) 的纵向分量。于是可以在一个使代数最简、但仍满足 \(\dot a^\mu\neq0\) 的一组局域 Taylor 数据上确定 \(C\)。

固定事件 \(t=0\)，取其瞬时静止系，并选择该事件满足
\[
 \vb v(0)=\vb0,
 \qquad
 \vb a(0)=\vb0,
 \qquad
 \vb j:=\left.\frac{\dd\vb a}{\dd t}\right|_{0}\neq\vb0.
\]
局部解析性、协变性、宇称和维数已经确定了唯一允许的张量结构；这组 Taylor 数据只用于计算其标量系数。

在该事件附近，源的位置、速度和加速度的 Taylor 展开分别为
\[
 \vb z(t')=\frac16\vb j\,t'^3+O(t'^4),
 \qquad
 \vb v(t')=\frac12\vb j\,t'^2+O(t'^3),
 \qquad
 \vb a(t')=\vb j\,t'+O(t'^2).
\]
取观测点 \(\vb x=\ell\vb n\)，\(|\vb n|=1\)。令 \(\epsilon=-1\) 表示推迟根、\(\epsilon=+1\) 表示超前根。光锥条件为
\[
 c^2t'^2
 =\left|\ell\vb n-\frac16\vb j\,t'^3+O(t'^4)\right|^2
 =\ell^2-\frac13\ell(\vb n\cdot\vb j)t'^3+O(\ell^5).
\]
作
\[
 t'_{\epsilon}=\epsilon\frac{\ell}{c}+B\ell^3+O(\ell^4)
\]
并代回，上式的 \(O(\ell^4)\) 项给出
\[
 2\epsilon cB\ell^4
 =-\epsilon\frac{\vb n\cdot\vb j}{3c^3}\ell^4,
\]
故
\begin{equation*}
 t'_{\epsilon}
 =\epsilon\frac{\ell}{c}
 -\frac{\vb n\cdot\vb j}{6c^4}\ell^3
 +O(\ell^4).
\end{equation*}
这里的 \(O(\ell^3)\) 根修正使 $R\cdot u$ 的相对 $O(\ell^2)$ 项变化；它与 $O(\ell^{-2})$ Coulomb 场相乘，留下有限贡献。

把正文\eqnrefs{eq:phys-rr-unified-ret-adv-field}中的固有时导数先显式完成，可写成
\[
 F_{\eta}^{\mu\nu}
 =\eta\frac{q}{4\pi\varepsilon_0c}
 \left[
 \frac{R^\mu a^\nu-R^\nu a^\mu}{(R\cdot u)^2}
 +\frac{c^2-R\cdot a}{(R\cdot u)^3}
 (R^\mu u^\nu-R^\nu u^\mu)
 \right]_{\tau_\eta},
\]
其中 \(\eta=+1\) 为推迟支，\(\eta=-1\) 为超前支，因此 \(\eta=-\epsilon\)。在当前瞬时静止系中，只需计算 \(F^{i0}\)。由上面的 Taylor 展开，
\[
\begin{aligned}
 R\cdot u
 &=-c^2t'-\frac12\ell(\vb n\cdot\vb j)t'^2+O(\ell^4),\\
 R\cdot a
 &=-\ell(\vb n\cdot\vb j)t'+O(\ell^3),
\end{aligned}
\]
并且
\[
 R^ia^0-R^0a^i
 =c j_i t'^2+O(\ell^3),
\]
\[
 R^iu^0-R^0u^i
 =c\ell n_i+\frac{c}{3}j_i t'^3+O(\ell^4).
\]
将光锥根 \(t'_{\epsilon}\) 代入并展开到 \(O(\ell^0)\)。第一项给出
\[
 \eta\frac1c
 \frac{R^ia^0-R^0a^i}{(R\cdot u)^2}
 =-\epsilon\frac{j_i}{c^4}+O(\ell),
\]
第二项给出
\[
 \eta\frac1c
 \frac{c^2-R\cdot a}{(R\cdot u)^3}
 (R^iu^0-R^0u^i)
 =\frac{n_i}{c\ell^2}
 +\epsilon\frac{j_i}{3c^4}+O(\ell).
\]
两项相加，因此
\[
 F_{\epsilon}^{i0}
 =\frac{q}{4\pi\varepsilon_0}
 \left[
 \frac{n_i}{c\ell^2}
 -\epsilon\frac{2j_i}{3c^4}
 \right]+O(\ell).
\]
推迟支对应 \(\epsilon=-1\)，超前支对应 \(\epsilon=+1\)，故 Coulomb 奇异项完全相同，而有限分支奇项为
\[
 \frac12\left(F_{\rm ret}^{i0}-F_{\rm adv}^{i0}\right)
 =\frac{q}{6\pi\varepsilon_0c^4}j_i.
\]
在此瞬时静止系，分支奇的有限磁场为零，于是
\[
 \frac12(\vb E_{\rm ret}-\vb E_{\rm adv})
 =\frac{q}{6\pi\varepsilon_0c^3}\vb j,
 \qquad
 \frac12(\vb B_{\rm ret}-\vb B_{\rm adv})=\vb0.
\]

最后恢复协变形式。在瞬时静止系中 \(u^\mu=(c,\vb0)\)，而当前选择的局域 Taylor 数据 有 \(a^\mu=0\)、\(\dot a^\mu=(0,\vb j)\)。唯一与上述分量一致的局域反对称张量是
\[
 F_{\rm R}^{\mu\nu}(z(\tau))
 =\frac{q}{6\pi\varepsilon_0c^5}
 \left(\dot a^\mu u^\nu-\dot a^\nu u^\mu\right).
\]
由于系数已在一个允许的局域 Taylor 数据 上唯一确定，且末式是 Lorentz 张量等式，它对一般 \(a^\mu\neq0\) 的世界线同样成立。这就是正文\eqnrefs{eq:phys-rr-regular-field-worldline}。

\end{proof}
